\documentclass[12pt]{article}
\usepackage{amsmath, amssymb, amsthm, geometry, setspace, graphicx,appendix,xcolor,tikz}
\usepackage{pgfplots}
\pgfplotsset{compat=1.18}
\usepackage[style=authoryear, hyperref=auto,dashed=false,maxnames=99,labelnumber=true,defernumbers=false,sorting=nyt,backend=biber]{biblatex}
\addbibresource{biblio.bib}
\newtheorem{proposition}{Proposition}
\usepackage{enumitem}
\setlist[itemize]{itemsep=0pt, topsep=2pt}
\geometry{margin=2.5cm}
\setstretch{1.2}

\begin{document}
\title{Intergenerational Theoretical Model with Formal and Informal Employment}
\maketitle 


\section{Introducción}

Desde la década de 1980, la desigualdad salarial ha aumentado de manera significativa tanto en países desarrollados como en economías en desarrollo. Una explicación central de este fenómeno es el cambio tecnológico sesgado hacia la habilidad (\emph{skill-biased technological change}), que ha incrementado la demanda relativa por trabajadores más calificados \parencite{Katz_1992,Krusell_2000,Card_2002,Buera_2021}. En esta línea, \textcite{Murphy_1993} argumentan que una fracción sustancial del aumento en la desigualdad salarial se explica por mayores retornos a la habilidad.

Este incremento en la desigualdad tiene implicancias directas para la movilidad intergeneracional. La evidencia empírica sintetizada en la denominada \emph{Great Gatsby Curve} muestra que economías con mayor persistencia intergeneracional tienden a exhibir mayores niveles de desigualdad \parencite{Corak_2013,BouchardSt-Amant2024,Durlauf_2018}. Este patrón sugiere que los procesos que determinan la distribución del ingreso dentro de una generación también influyen en la transmisión de ventajas económicas entre generaciones, planteando la pregunta sobre los mecanismos que vinculan desigualdad y movilidad social.

Un canal natural para entender esta relación es la inversión en capital humano. En un contexto en el que el mercado valora crecientemente las habilidades, resulta plausible que los padres enfrenten mayores incentivos a invertir en el capital humano de sus hijos con el objetivo de mejorar su bienestar económico futuro. Bajo retornos homogéneos al capital humano, un aumento en los retornos a la habilidad induciría incrementos relativamente uniformes en la inversión parental, desplazando la distribución de ingresos de los hijos sin alterar sustancialmente su composición relativa.

Sin embargo, cuando los retornos de mercado a las habilidades dependen del nivel de capital humano alcanzado, y son crecientes en dicho nivel,  las decisiones de inversión parental pueden realizar cambios discretos y considerables frente a diferencias marginales en el ingreso del padre. Esta estructura es análoga a un sistema impositivo con tasas marginales que varían por tramos: así como los incentivos a generar ingreso cambian al cruzar determinados umbrales de ingreso, los incentivos a invertir en capital humano se modifican al superar niveles específicos de capital humano potencial. A diferencia del caso impositivo, donde las tasas marginales suelen reducir los retornos netos, en este contexto los retornos marginales aumentan al cruzar dichos umbrales, amplificando las diferencias en los incentivos a invertir.

Como consecuencia, los padres con mayores recursos pueden financiar inversiones suficientes para alcanzar los niveles de capital humano más valorados por el mercado y beneficiarse de retornos crecientes, mientras que los padres con menores ingresos, al no poder cubrir el costo necesario para superar estos umbrales, terminan concentrándose en niveles bajos de acumulación de capital humano. Este mecanismo genera disparidades persistentes en la inversión parental y sienta las bases para una menor movilidad intergeneracional en los tramos mas altos de ingresos.

Los retornos al capital humano pueden volverse crecientes a partir de ciertos umbrales por diversos mecanismos. Por un lado, el surgimiento de economías de estrellas \parencite{Rosen_1983,Autor_2020} ha incrementado la demanda por trabajadores altamente calificados en firmas intensivas en talento, elevando desproporcionadamente los retornos en los tramos superiores de la distribución de capital humano. En este contexto, alcanzar niveles suficientemente altos de capital humano permite acceder a ocupaciones y firmas donde las habilidades acumuladas son especialmente valoradas y, en consecuencia, mejor remuneradas.
Por otro lado, la acumulación de capital humano por encima de ciertos niveles facilita el acceso a colegios y universidades de élite, así como a empleadores de firmas más productivas que valoran de manera diferencial dichas credenciales \parencite{Zimmerman_2019,Chetty_2025}. El acceso a estos entornos no solo incrementa directamente los retornos al capital humano invertido, sino que también habilita la formación de redes y conexiones relevantes \parencite{Bavaro_2022}, lo que puede amplificar disparidades iniciales y generar ineficiencias asociadas a una mala asignación del talento \parencite{Bandyopadhyay_2019}.


En este trabajo seguimos el marco clásico de \textcite{Becker_Tomes_1979}, con simplificaciones derivadas de \textcite{Solon_1999}, incorporando retornos diferenciados a la inversión en capital humano. Esta extensión modifica las decisiones óptimas de inversión parental y, en consecuencia, las propiedades de la movilidad intergeneracional de ingresos.

Como resultado, el problema de inversión parental presenta no convexidades asociadas a la presencia de regímenes diferenciados de retornos al capital humano. En particular, el padre enfrenta un régimen de bajos retornos para niveles reducidos de inversión y un régimen de altos retornos una vez que se supera un determinado umbral. En este contexto, existe un nivel de ingreso parental para el cual el padre es indiferente entre realizar una inversión que lo mantiene en el régimen de bajos retornos y una inversión que le permite acceder al régimen de altos retornos. Dado este punto de indiferencia, los niveles intermedios de inversión entre ambos regímenes resultan estrictamente dominados y no son elegidos en equilibrio. Esta estructura implica un punto de bifurcación en el cual pequeñas variaciones en el ingreso del padre pueden inducir ajustes discretos y de gran magnitud en la inversión en capital humano del hijo.

Este mecanismo es conceptualmente análogo al que surge en modelos clásicos de impuestos no lineales \parencite{Burtless_1978,Hausman_1985}, como el programa de impuesto negativo a la renta (NIT). En dichos modelos, la introducción de transferencias con tasas marginales discontinuas genera conjuntos presupuestarios no convexos, de modo que existen intervalos de oferta laboral que nunca contienen un óptimo global bajo curvas de indiferencia convexas.\footnote{Esto difiere de los modelos de \emph{bunching} \parencite{Saez_2010}, en los que impuestos progresivos con tasas marginales crecientes generan conjuntos presupuestarios convexos, los cuales poseen una solución global única que puede ubicarse en el kink de tasas marginales. En contraste, en conjuntos presupuestarios no convexos el kink no constituye una solución global bajo ninguna configuración de parámetros, asumiendo curvas de indiferencia convexas.} 
De manera similar, en nuestro marco la presencia de umbrales en los retornos al capital humano implica que ciertos niveles intermedios de inversión parental están estrictamente dominados, concentrando las decisiones en regímenes de baja o alta inversión. Este patrón contribuye a la formación de dinastías económicas y a una reducción de la movilidad intergeneracional, particularmente en los tramos altos de la distribución del ingreso.

El modelo contribuye a la comprensión de las no linealidades en la elasticidad intergeneracional de ingresos (IGE), caracterizadas por una mayor persistencia en los extremos de la distribución. En particular, los hogares de mayores ingresos tienden a exhibir una IGE más elevada, reflejando una menor movilidad intergeneracional. Esta predicción se vincula con una literatura empírica que documenta heterogeneidad en la movilidad a lo largo de la distribución del ingreso parental. \textcite{Black_2011} presenta evidencia de no linealidades en la IGE y discute mecanismos como restricciones crediticias \parencite{Lochner_2011,Lochner_2020}, que resultan clave para comprender la baja movilidad intergeneracional en los sectores de menores ingresos. Por su parte, \textcite{Gregg_2019} documenta empíricamente una relación en forma de ``J'' de la IGE, en la que los deciles más bajos y más altos de la distribución presentan una menor movilidad intergeneracional. Finalmente, \textcite{Durlauf_1996,Durlauf_2018} enfatizan el rol de los efectos de vecindario y la segregación residencial como determinantes de la acumulación de capital humano, reforzando la persistencia intergeneracional en los polos de la distribución.

Nuestro enfoque también se relaciona con \textcite{Becker_2018}, quienes muestran que el ingreso de los hijos es convexo respecto al capital humano parental y que pueden emerger dinastías económicas. En su marco, este resultado surge de complementariedades en la producción de capital humano, en el sentido de que el capital humano de los padres incrementa la productividad tanto en el mercado laboral como en la formación de habilidades de los hijos. En contraste, en nuestro modelo la convexidad del ingreso y la persistencia intergeneracional emergen de diferencias en los retornos al capital humano a partir de ciertos niveles de ingreso, incluso en ausencia de complementariedades tecnológicas en la producción de habilidades.

Finalmente, el trabajo se relaciona con \textcite{Solon_2004}, quien muestra que aumentos globales en los retornos al capital humano, combinados con mayores niveles de desigualdad, incrementan la elasticidad intergeneracional del ingreso, en línea con la Great Gatsby Curve. Sin embargo, en dicho enfoque el aumento de la desigualdad es exógeno y no se evalúa el rol de las decisiones de inversión parental ni el impacto de cambios endógenos en la estructura de retornos del capital humano sobre la transmisión intergeneracional de la desigualdad.



%\textcite{Galor_1993}
%Rentas y college : \textcite{Zimmerman_2019,Murphy_1991,Torvik_2002}
%Formación de capital Humano \textcite{Cunha_2007}
%Articulos de Chetty:  \textcite{Chetty_2014,Chetty_2025}
%Credit Constrains: \textcite{Lochner_2011,Lochner_2020}
%Referals \textcite{Bavaro_2022}
%Missallocation: \textcite{Bandyopadhyay_2019}



\section{Modelo}

Considere una familia con dos generaciones. El padre vive en el período $t-1$ y decide cómo asignar su ingreso entre consumo propio e inversión en el capital humano de su hijo, quien vive en el período $t$. Una mayor inversión en capital humano hoy se traduce en mayores ingresos del hijo en el futuro.

En el caso más simple, el ingreso del hijo aumenta de manera proporcional con la inversión del padre: cada unidad adicional invertida genera el mismo retorno marginal, independientemente del nivel de inversión. En este entorno, cambios en los retornos a la inversión en capital humano afectan de manera similar a todas las familias y no alteran cualitativamente las decisiones óptimas entre agentes. Este trabajo, en cambio, explora un mecanismo distinto. Supóngase que los retornos a la inversión en capital humano no son constantes, sino que aumentan a partir de cierto umbral de inversión. En particular, a partir de un nivel $\bar I$, cada unidad adicional invertida genera un retorno marginal mayor que el previo.

Esta estructura introduce una no linealidad en la transmisión intergeneracional de ingresos. El padre maximiza una función de utilidad que es una suma ponderada de dos componentes concavos: la utilidad derivada del consumo propio y la utilidad asociada al ingreso futuro del hijo. Bajo estas preferencias, las curvas de indiferencia son convexas.

Como se ilustra en la Figura~\ref{fig:kink}, la presencia de un \emph{kink} en los retornos de mercado del capital humano puede generar comportamientos no continuos y cambios discretos en los niveles óptimos de inversión parental. Para un nivel específico de ingreso del padre, pueden coexistir dos óptimos globales de inversión: uno localizado en el régimen de bajos retornos y otro que supera el umbral $\bar I$ y se beneficia del aumento discreto en los retornos marginales. Un incremento marginal en el ingreso parental hace que el óptimo global se desplace al régimen de altos retornos, mientras que una reducción marginal en dicho ingreso lleva al padre a elegir el régimen de bajos retornos. Cabe destacar que, bajo preferencias concavas en consumo e ingreso del hijo, el padre nunca elige invertir exactamente en $\bar I$, ya que dicha inversión no permite aprovechar los retornos marginales crecientes y resulta estrictamente dominada.


\begin{tikzpicture}[xscale=1.3, yscale=0.5]
	\label{fig:kink}
	% Parameters
	%------------------------------------------------
	\def\rone{-0.60}   % very low first slope
	\def\rtwo{1.80}    % very high second slope
	\def\Ibar{4}
	\def\e{1}
	
	\def\IL{2}
	\def\IH{6}
	
	% Technology
	\def\yone(#1){(1+\rone)*(#1)+\e}
	\def\ytwo(#1){(1+\rtwo)*(#1)+\e-(\rtwo-\rone)*\Ibar}
	
	%------------------------------------------------
	% Axes
	%------------------------------------------------
	\draw[->] (0,0) -- (8.5,0) node[below] {$I_{t-1}$};
	\draw[->] (0,0) -- (0,13) node[left] {$y_t$};
	
	%------------------------------------------------
	% Child income (technology)
	%------------------------------------------------
	\draw[thick,blue]
	plot[domain=0:\Ibar] (\x,{\yone(\x)});
	\draw[thick,blue]
	plot[domain=\Ibar:7.8] (\x,{\ytwo(\x)});
	
	
	%------------------------------------------------
	% Kink
	%------------------------------------------------
	\filldraw[black] (\Ibar,{\yone(\Ibar)}) circle (2pt);
	\draw[dashed] (\Ibar,0) -- (\Ibar,{\yone(\Ibar)});
	\node[below] at (\Ibar,0) {$\bar I$};
	
	%------------------------------------------------
	% Double tangency points
	%------------------------------------------------
	\filldraw[red] (\IL,{\yone(\IL)}) circle (2pt);
	\node[below left] at (\IL,{\yone(\IL)}) {$I_L$};
	
	\filldraw[red] (\IH,{\ytwo(\IH)}) circle (2pt);
	\node[below right] at (\IH,{\ytwo(\IH)}) {$I_H$};
	%------------------------------------------------
	% Convex indifference curve (lowered but still convex)
	%------------------------------------------------
	\draw[thick]
	plot[smooth,domain=0.7:7.6]
	(\x,{
		0.3*(\x-\IL)^2
		+ (1+\rone)*(\x-\IL)
		+ \yone(\IL)
	});
	
	\node at (3,1) {$(1+r)$};
	\node at (5.5,4) {$(1+r')$};
	%------------------------------------------------
	% Tangency guides
	%------------------------------------------------
	% At I_L
	\draw[dotted]
	(\IL-1.1,{\yone(\IL)-(1+\rone)})
	--
	(\IL+1.1,{\yone(\IL)+(1+\rone)});
	
	% At I_H
	\draw[dotted]
	(\IH-1.1,{\ytwo(\IH)-(1+\rtwo)})
	--
	(\IH+1.1,{\ytwo(\IH)+(1+\rtwo)});
	
\end{tikzpicture}

\section{Implicaciones para la acumulación de capital humano}

Considere una familia con dos generaciones. El padre vive en el periodo $t-1$ y decide cómo asignar su ingreso $y_{t-1}$ entre consumo propio e inversión en el capital humano de su hijo $I_{t-1|}$, quien vive en el periodo $t$. 
Tenemos en cuenta que el ingreso del hijo presenta puntos en los cuales el precio del capital humano cambia de manera discreta desde una tasa $r$,  a una tasa  $r'$ a partir de una inversion de capital humano $\bar{I}$. En este caso $e_t$ son componentes heredados del capital humano.

\[
y_{t} = (1+r)I_{t-1}+e_t+(r'-r)\max\{0,I_{t+1}- \bar{I}\}
\]

Siguiendo el modelo simplificado de \textcite{Solon_1999}, el padre maximiza su utilidad en consumo y en el ingreso del hijo, ponderados por un factor de relevancia 

$$\alpha \log (c_{t-1}) + (1-\alpha)\log y_t$$

Como mostramos en la seccion anterior, el proceso de cambios distretos en los retornos del capital humano en conjunto con utilidades concavas en el consumo e ingreso del hijo , implican curvas de indiferencia convexas. y una situacion en la que para los padres no es optimo invertir en el punto (y su vecindad) donde se genera el cambio de paradigma en las tasas de retornos. Este elemento introduce una no linealidad en la tecnología de transmisión intergeneracional de ingresos. 

Las soluciones locales del problema del padre abogan a obtener la inversion optima de cada segmento de la relacion lineal entre los ingresos del hijo y la inversion en capital humano, asumiendo que dicha relacion se mantiene por todo el dominio de inversion.  
Obviamente para ser un candidato a solucion global, la inversion optima de cada segmento debe estar en el dominio del segmento manifestado en la relacion lineal de ingresos. Una vez se tienen los candidatos a optimo global, se comparan las funciones de utilidad indirecta de cada candidato. 

Aunque es mas util seguir el proceso de encontrar el punto de ingresos del padre en el que el agente esta indiferente entre invertir en el regimen bajo y en el regimen alto. 


Dado un régimen tecnológico con retorno constante $\rho\in\{r,r'\}$, el padre resuelve el siguiente problema local:
\[
\max_{I\in[0,y_{t-1}]}
\;
\alpha \log(y_{t-1}-I)
+
(1-\alpha)\log\big((1+\rho)I+e_t\big).
\]

La condición de primer orden está dada por
\[
\frac{\alpha}{y_{t-1}-I}
=
\frac{(1-\alpha)(1+\rho)}{(1+\rho)I+e_t},
\]
de donde se obtiene la inversión óptima local
\[
I^*(\rho)
=
(1-\alpha)y_{t-1}
-
\frac{\alpha}{1+\rho}e_t.
\]

Evaluando la función objetivo en $I^*(\rho)$ y utilizando la condición de primer orden, la utilidad indirecta asociada al régimen $\rho$ puede escribirse como
\[
V(\rho)
=
\log\big((1+\rho)y_{t-1}+e_t\big)
-
\alpha\log(1+\rho)
+
\alpha\log\alpha
+
(1-\alpha)\log(1-\alpha).
\]
El último término es constante en $\rho$ y, por lo tanto, irrelevante para la comparación entre regímenes.


El padre es indiferente entre invertir bajo el régimen de bajo retorno $r$ y el de alto retorno $r'>r$ cuando
\[
V(r)=V(r').
\]
Utilizando la expresión cerrada para la utilidad indirecta, esta condición puede escribirse como
\[
\log\frac{(1+r')y_{t-1}+e_t}{(1+r)y_{t-1}+e_t}
=
\alpha
\log\frac{1+r'}{1+r}.
\tag{IC}
\]

Exponenciando ambos lados de \eqref{IC} se obtiene
\[
\frac{(1+r')y_{t-1}+e_t}{(1+r)y_{t-1}+e_t}
=
\left(\frac{1+r'}{1+r}\right)^{\alpha}.
\]

Multiplicando y reordenando términos:
\[
(1+r')y_{t-1}+e_t
=
\left(\frac{1+r'}{1+r}\right)^{\alpha}
\big[(1+r)y_{t-1}+e_t\big].
\]

Agrupando los términos en $y_{t-1}$:
\[
\Big[(1+r')-(1+r)\Big(\tfrac{1+r'}{1+r}\Big)^{\alpha}\Big]y_{t-1}
=
e_t\Big[\Big(\tfrac{1+r'}{1+r}\Big)^{\alpha}-1\Big].
\]

Por lo tanto, el nivel de ingreso del padre que lo deja indiferente entre ambos regímenes está dado por
\[
\boxed{
y_{t-1}^{\text{ind}}(e_t)
=
\frac{
e_t\Big[\big(\tfrac{1+r'}{1+r}\big)^{\alpha}-1\Big]
}{
(1+r')-(1+r)\big(\tfrac{1+r'}{1+r}\big)^{\alpha}
}
}
\]

Dado que $r'>r$ y $\alpha\in(0,1)$, el numerador y el denominador de la expresión anterior son estrictamente positivos, lo que garantiza que $y_{t-1}^{\text{ind}}(e_t)>0$. En consecuencia, para $y_{t-1}<y_{t-1}^{\text{ind}}(e_t)$ el padre prefiere invertir bajo el régimen de bajo retorno, mientras que para $y_{t-1}>y_{t-1}^{\text{ind}}(e_t)$ resulta óptimo saltar discretamente al régimen de alto retorno. Este mecanismo introduce una discontinuidad en la inversión óptima y genera una no linealidad en la tecnología de transmisión intergeneracional de ingresos.


Para que este nivel de ingreso actúe efectivamente como un umbral (\emph{binding}) de cambio de régimen, es necesario que, evaluado bajo el régimen de bajo retorno, la elección óptima de inversión alcance el punto de quiebre de la tecnología. En particular, el ingreso \(y_{t-1}^{\text{ind}}(e_t)\) debe ser lo suficientemente bajo como para que la inversión óptima en el primer segmento sea igual a \(\bar I\).

Recordando que, bajo el régimen de bajo retorno, la inversión óptima viene dada por
\[
I^*(y_{t-1},e_t)
=
(1-\alpha)y_{t-1}-\frac{\alpha}{1+r}e_t,
\]
la condición para que el umbral \(y_{t-1}^{\text{ind}}(e_t)\) sea económicamente relevante es
\[
I^*\big(y_{t-1}^{\text{ind}}(e_t),e_t\big)\;\leq\;\bar I.
\]

Cuando esta condición se satisface, el padre se encuentra exactamente en el límite del primer segmento de la tecnología y el cruce en preferencias entre regímenes se traduce en un cambio efectivo de régimen de inversión. En caso contrario, si para \(y_{t-1}^{\text{ind}}(e_t)\) la inversión óptima bajo el régimen de bajo retorno excede \(\bar I\), el punto de indiferencia carece de contenido económico, ya que la restricción tecnológica no es vinculante y el padre ya estaría operando en el segmento de alto retorno.


Con lo anterior, la tecnología de transmisión intergeneracional de ingresos puede escribirse en forma reducida reemplazando las inversiones óptimas en la ecuación de ingreso del hijo. Dado que la elección de régimen depende del nivel de ingreso del padre, se obtienen dos relaciones intergeneracionales distintas.

\paragraph{Régimen de bajo retorno.}
Para padres con ingresos por debajo del umbral de indiferencia,
\(y_{t-1} \leq y_{t-1}^{\text{ind}}(e_t)\),
la inversión óptima se realiza bajo el régimen de bajo retorno \(r\),
\[
I^*(r)
=
(1-\alpha)y_{t-1}
-
\frac{\alpha}{1+r}e_t.
\]
Sustituyendo en la ecuación de ingreso del hijo se obtiene
\[
\begin{aligned}
y_t
&=
(1+r)I^*(r)+e_t \\
&=
(1+r)(1-\alpha)y_{t-1}
+
\alpha e_t.
\end{aligned}
\]

\paragraph{Régimen de alto retorno.}
Para padres con ingresos superiores al umbral,
\(y_{t-1} > y_{t-1}^{\text{ind}}(e_t)\),
el óptimo global implica invertir bajo el régimen de alto retorno \(r'\),
\[
I^*(r')
=
(1-\alpha)y_{t-1}
-
\frac{\alpha}{1+r'}e_t.
\]
En este caso, el ingreso del hijo viene dado por
\[
\begin{aligned}
y_t
&=
(1+r')I^*(r')+e_t \\
&=
(1+r')(1-\alpha)y_{t-1}
+
\alpha e_t.
\end{aligned}
\]

En consecuencia, la relación intergeneracional de ingresos adopta una forma no lineal y por tramos:
\[
y_t
=
\begin{cases}
(1+r)(1-\alpha)y_{t-1}+\alpha e_t,
& \text{si } y_{t-1}\leq y_{t-1}^{\text{ind}}(e_t), \\[0.4em]
(1+r')(1-\alpha)y_{t-1}+\alpha e_t,
& \text{si } y_{t-1}> y_{t-1}^{\text{ind}}(e_t).
\end{cases}
\]

Esta expresión muestra que un cambio discreto en los retornos del capital humano se traduce en un quiebre en la pendiente de la curva de transmisión intergeneracional de ingresos, generando heterogeneidad en la persistencia intergeneracional según el nivel de ingreso del padre.

\section{Implicaciones para la acumulación de capital humano}

Considere una familia con dos generaciones. El padre vive en el periodo $t-1$ y decide cómo asignar su ingreso $y_{t-1}$ entre consumo propio e inversión en el capital humano de su hijo $I_{t-1}$, quien vive en el periodo $t$. El ingreso del hijo depende de la inversión parental y de una dotación heredada $e_t$, y presenta un cambio discreto en los retornos del capital humano una vez que el capital humano efectivo supera un umbral $h>0$.

La tecnología de generación de ingresos del hijo está dada por
\[
y_t
=
(1+r)I_{t-1}
+
e_t
+
r'\max\{0,I_{t-1}+e_t-h\},
\]
donde $r'>0$ representa un retorno adicional que se aplica únicamente al excedente del capital humano efectivo por sobre el umbral $h$.

Definiendo el umbral efectivo de inversión
\[
\bar I(e_t)\equiv h-e_t,
\]
la tecnología puede escribirse de manera equivalente por tramos como
\[
y_t =
\begin{cases}
(1+r)I_{t-1}+e_t,
& \text{si } I_{t-1}\le \bar I(e_t), \\[0.4em]
(1+r+r')I_{t-1}+(1+r')e_t,
& \text{si } I_{t-1}> \bar I(e_t).
\end{cases}
\]

Siguiendo el modelo simplificado de \textcite{Solon_1999}, el padre maximiza una función de utilidad logarítmica en consumo e ingreso del hijo,
\[
\alpha \log(c_{t-1})+(1-\alpha)\log y_t,
\]
donde $c_{t-1}=y_{t-1}-I_{t-1}$ y $\alpha\in(0,1)$.

Debido a la concavidad de la utilidad y a la presencia de un quiebre en la tecnología, el problema del padre puede resolverse comparando soluciones locales en cada régimen tecnológico.

\subsection*{Óptimos locales}

\paragraph{Régimen de bajo retorno.}
Cuando $I_{t-1}\le \bar I(e_t)$, el padre resuelve
\[
\max_{I\in[0,y_{t-1}]}
\;
\alpha\log(y_{t-1}-I)
+
(1-\alpha)\log\big((1+r)I+e_t\big).
\]
La condición de primer orden es
\[
\frac{\alpha}{y_{t-1}-I}
=
\frac{(1-\alpha)(1+r)}{(1+r)I+e_t},
\]
de donde se obtiene la inversión óptima local
\[
I_L^*
=
(1-\alpha)y_{t-1}
-
\frac{\alpha}{1+r}e_t.
\]

\paragraph{Régimen de alto retorno.}
Cuando $I_{t-1}>\bar I(e_t)$, el problema local es
\[
\max_{I\in[0,y_{t-1}]}
\;
\alpha\log(y_{t-1}-I)
+
(1-\alpha)\log\big((1+r+r')I+(1+r')e_t\big).
\]
La condición de primer orden es
\[
\frac{\alpha}{y_{t-1}-I}
=
\frac{(1-\alpha)(1+r+r')}{(1+r+r')I+(1+r')e_t},
\]
lo que implica una inversión óptima local
\[
I_H^*
=
(1-\alpha)y_{t-1}
-
\frac{\alpha(1+r')}{1+r+r'}e_t.
\]

\subsection*{Comparación entre regímenes y umbral de indiferencia}

Evaluando la utilidad en las soluciones óptimas y utilizando las condiciones de primer orden, las utilidades indirectas asociadas a cada régimen pueden escribirse como
\[
V_L
=
\log\big((1+r)y_{t-1}+e_t\big)
-
\alpha\log(1+r)
+
C,
\]
\[
V_H
=
\log\big((1+r+r')y_{t-1}+(1+r')e_t\big)
-
\alpha\log(1+r+r')
+
C,
\]
donde $C=\alpha\log\alpha+(1-\alpha)\log(1-\alpha)$ es constante e irrelevante para la comparación.

El padre es indiferente entre invertir bajo el régimen de bajo retorno y el de alto retorno cuando $V_L=V_H$, lo que implica
\[
\log\frac{(1+r+r')y_{t-1}+(1+r')e_t}{(1+r)y_{t-1}+e_t}
=
\alpha\log\frac{1+r+r'}{1+r}.
\]

Resolviendo esta expresión se obtiene el nivel de ingreso del padre que lo deja indiferente entre ambos regímenes:
\[
\boxed{
y_{t-1}^{\text{ind}}(e_t)
=
\frac{
e_t\Big[\big(\tfrac{1+r+r'}{1+r}\big)^{\alpha}-(1+r')\Big]
}{
(1+r+r')-(1+r)\big(\tfrac{1+r+r'}{1+r}\big)^{\alpha}
}
}
\]

Para que este nivel de ingreso actúe efectivamente como un umbral de cambio de régimen, es necesario que la inversión óptima bajo el régimen de bajo retorno no supere el punto de quiebre tecnológico. En particular, debe cumplirse
\[
I_L^*\big(y_{t-1}^{\text{ind}}(e_t),e_t\big)
\le
\bar I(e_t)=h-e_t.
\]

\subsection*{Relación intergeneracional de ingresos}

Sustituyendo las inversiones óptimas en la tecnología de ingresos del hijo, se obtiene la siguiente relación intergeneracional por tramos:
\[
y_t
=
\begin{cases}
(1+r)(1-\alpha)y_{t-1}+\alpha e_t,
& \text{si } y_{t-1}\le y_{t-1}^{\text{ind}}(e_t), \\[0.4em]
(1+r+r')(1-\alpha)y_{t-1}
+\alpha(1+r')e_t,
& \text{si } y_{t-1}> y_{t-1}^{\text{ind}}(e_t).
\end{cases}
\]

Esta expresión muestra que la transmisión intergeneracional de ingresos es no lineal y dependiente del historial: familias que parten con ingresos suficientemente altos acceden a un régimen de mayores retornos del capital humano, lo que amplifica la persistencia intergeneracional y genera heterogeneidad endógena en la movilidad económica.


\subsection{Dinámica de largo plazo y dependencia del historial}

Considérese dos dinastías, $i$ y $j$, idénticas en preferencias y tecnología, que difieren únicamente en sus condiciones iniciales. Denótese por
\[
y^{\text{ind}}(\bar e)
\]
el nivel de ingreso parental que deja indiferente al padre entre invertir en el régimen de bajo y alto retorno cuando la dotación heredada es constante e igual a $\bar e$.

Las condiciones iniciales de ambas dinastías están dadas por
\[
y_0^i=\frac{3}{2}y^{\text{ind}}(\bar e),
\qquad
y_0^j=\frac{1}{2}y^{\text{ind}}(\bar e),
\]
de modo que la dinastía $i$ parte estrictamente por sobre el umbral de indiferencia, mientras que la dinastía $j$ parte por debajo.

---

\subsubsection*{Proceso de dotación heredada}

La dotación heredada evoluciona de manera distinta para cada dinastía. Para la dinastía rica $i$, se asume una dotación constante en el tiempo:
\[
e_t^i=\bar e \qquad \forall t.
\]

En contraste, la dinastía pobre $j$ experimenta un crecimiento determinístico de su dotación heredada a una tasa $g>0$:
\[
e_t^j=(1+g)^t\bar e.
\]

Este supuesto captura la idea de que las familias inicialmente desfavorecidas pueden experimentar mejoras graduales en su entorno, redes o capital social, aun cuando partan de una posición inicial baja.

---

\subsubsection*{Trayectorias intergeneracionales}

La relación intergeneracional de ingresos viene dada por
\[
y_t^k
=
\begin{cases}
\beta_L y_{t-1}^k + \alpha e_t^k,
& \text{si } y_{t-1}^k\le y^{\text{ind}}(e_t^k), \\[0.4em]
\beta_H y_{t-1}^k + \alpha(1+r')e_t^k,
& \text{si } y_{t-1}^k> y^{\text{ind}}(e_t^k),
\end{cases}
\quad k\in\{i,j\},
\]
donde
\[
\beta_L=(1+r)(1-\alpha),
\qquad
\beta_H=(1+r+r')(1-\alpha),
\qquad
\beta_H>\beta_L.
\]

---

\paragraph{Dinastía rica.}
Dado que $y_0^i>y^{\text{ind}}(\bar e)$ y $e_t^i=\bar e$ para todo $t$, la dinastía $i$ permanece en el régimen de alto retorno en todos los periodos. Su trayectoria satisface
\[
y_t^i
=
\beta_H y_{t-1}^i + \alpha(1+r')\bar e,
\]
lo que implica
\[
y_t^i
=
\beta_H^t y_0^i
+
\alpha(1+r')\bar e
\sum_{s=0}^{t-1}\beta_H^s.
\]

Si $\beta_H<1$, la dinastía converge a un estado estacionario alto; si $\beta_H\ge1$, su ingreso crece sin cota, amplificando permanentemente la ventaja inicial.

---

\paragraph{Dinastía pobre.}
La dinastía $j$ parte en el régimen de bajo retorno y su dinámica inicial está dada por
\[
y_t^j
=
\beta_L y_{t-1}^j + \alpha e_t^j
=
\beta_L y_{t-1}^j + \alpha (1+g)^t\bar e.
\]

Iterando hacia adelante se obtiene
\[
y_t^j
=
\beta_L^t y_0^j
+
\alpha\bar e
\sum_{s=0}^{t-1}\beta_L^s(1+g)^{t-s}.
\]

Mientras $y_t^j\le y^{\text{ind}}(e_t^j)$, la dinastía permanece atrapada en el régimen de bajo retorno. El cruce al régimen alto ocurre, si es que ocurre, en el primer periodo $\tau$ tal que
\[
y_\tau^j>y^{\text{ind}}(e_\tau^j).
\]

---

\subsubsection*{Comparación de largo plazo}

Aun cuando la dinastía pobre experimenta un crecimiento exógeno de su dotación heredada, su capacidad de capitalizar dicho crecimiento es limitada mientras permanece en el régimen de bajo retorno. En contraste, la dinastía rica amplifica tanto su ventaja inicial como cualquier nivel dado de dotación heredada debido a la mayor pendiente de la relación intergeneracional.

En particular, si el crecimiento de la dotación heredada satisface
\[
1+g \le \beta_H,
\]
entonces existe un conjunto no trivial de parámetros para el cual la dinastía pobre nunca alcanza el umbral de indiferencia, generándose una trampa de bajo ingreso intergeneracional. Por el contrario, solo si el crecimiento de la dotación heredada es suficientemente elevado, de modo que
\[
1+g>\beta_L,
\]
la dinastía pobre puede eventualmente cruzar el umbral y transitar hacia el régimen de alto retorno.

---

\subsubsection*{Interpretación}

Este ejercicio muestra que la convergencia intergeneracional no depende únicamente de mejoras graduales en las dotaciones heredadas, sino críticamente del régimen tecnológico en el que opera cada familia. Condiciones iniciales desfavorables pueden persistir en el tiempo incluso cuando la dotación heredada crece, debido a la menor capacidad de acumulación intergeneracional en el régimen de bajo retorno. En consecuencia, la historia inicial de cada dinastía determina de manera endógena su trayectoria de largo plazo.


\subsection{Dinámica local alrededor del umbral de indiferencia}

Considérese nuevamente la relación intergeneracional de ingresos por tramos:
\[
y_t
=
\begin{cases}
\beta_L y_{t-1}+\alpha e_t,
& \text{si } y_{t-1}\le y^{\text{ind}}(e_t), \\[0.4em]
\beta_H y_{t-1}+\alpha(1+r')e_t,
& \text{si } y_{t-1}> y^{\text{ind}}(e_t),
\end{cases}
\]
donde
\[
\beta_L=(1+r)(1-\alpha), 
\qquad 
\beta_H=(1+r+r')(1-\alpha),
\qquad 
\beta_H>\beta_L.
\]

Sea $\bar e$ una dotación heredada constante y defínase $y^{\text{ind}}(\bar e)$ como el nivel de ingreso parental que deja indiferente al agente entre invertir en el régimen de bajo y alto retorno.

---

\subsubsection*{Condiciones iniciales}

Considérense dos dinastías, $i$ y $j$, idénticas en todo salvo por una diferencia arbitrariamente pequeña en sus condiciones iniciales:
\[
y_0^i
=
y^{\text{ind}}(\bar e),
\qquad
y_0^j
=
y^{\text{ind}}(\bar e)-\varepsilon,
\qquad
\varepsilon>0 \text{ pequeño}.
\]

Ambas enfrentan inicialmente la misma dotación heredada:
\[
e_0^i=e_0^j=\bar e.
\]

Por construcción, la dinastía $i$ opera en el régimen de alto retorno desde el periodo inicial, mientras que la dinastía $j$ parte en el régimen de bajo retorno.

---

\subsubsection*{Trayectorias determinísticas sin shocks}

En ausencia de shocks o crecimiento adicional en la dotación heredada, las trayectorias vienen dadas por:

\paragraph{Dinastía en el régimen alto.}
\[
y_t^i
=
\beta_H y_{t-1}^i + \alpha(1+r')\bar e,
\]
con solución cerrada:
\[
y_t^i
=
\beta_H^t y^{\text{ind}}(\bar e)
+
\alpha(1+r')\bar e
\sum_{s=0}^{t-1}\beta_H^s.
\]

\paragraph{Dinastía en el régimen bajo.}
\[
y_t^j
=
\beta_L y_{t-1}^j + \alpha\bar e,
\]
lo que implica:
\[
y_t^j
=
\beta_L^t\big(y^{\text{ind}}(\bar e)-\varepsilon\big)
+
\alpha\bar e
\sum_{s=0}^{t-1}\beta_L^s.
\]

Dado que $\beta_H>\beta_L$, la brecha entre ambas dinastías crece en el tiempo aun cuando $\varepsilon$ sea arbitrariamente pequeño.

---

\subsubsection*{Crecimiento compensatorio de la dotación heredada}

Supóngase ahora que la dinastía $j$ experimenta un crecimiento exógeno de su dotación heredada a una tasa constante $g>0$, mientras que la dinastía $i$ mantiene $e_t^i=\bar e$:
\[
e_t^j=(1+g)^t\bar e.
\]

La dinámica de la dinastía $j$ pasa a ser:
\[
y_t^j
=
\beta_L y_{t-1}^j + \alpha(1+g)^t\bar e,
\]
con solución:
\[
y_t^j
=
\beta_L^t\big(y^{\text{ind}}(\bar e)-\varepsilon\big)
+
\alpha\bar e
\sum_{s=0}^{t-1}\beta_L^s(1+g)^{t-s}.
\]

La dinastía $j$ alcanza la trayectoria de la dinastía $i$ en el primer periodo $\tau$ tal que:
\[
y_\tau^j
\ge
y_\tau^i.
\]

---

\subsubsection*{Condición mínima de crecimiento}

Una condición necesaria para que la dinastía inicialmente desfavorecida pueda cerrar la brecha es que el crecimiento de la dotación heredada compense la menor pendiente intergeneracional:
\[
1+g \ge \frac{\beta_H}{\beta_L}.
\]

Si esta condición no se cumple, incluso una ventaja inicial infinitesimal para la dinastía $i$ genera divergencia permanente. En cambio, si
\[
1+g > \frac{\beta_H}{\beta_L},
\]
entonces existe un horizonte finito $\tau(\varepsilon)$ tal que la dinastía $j$ no solo cruza el umbral de indiferencia, sino que alcanza y eventualmente supera la trayectoria de la dinastía inicialmente favorecida.

---

\subsubsection*{Interpretación}

Este ejercicio muestra que la dinámica intergeneracional es localmente inestable alrededor del umbral de indiferencia. Diferencias iniciales arbitrariamente pequeñas pueden generar trayectorias persistentemente divergentes, a menos que la dinastía inicialmente desfavorecida experimente un crecimiento suficientemente alto de su dotación heredada. En este sentido, la movilidad intergeneracional depende críticamente no solo de la magnitud de los shocks o del crecimiento futuro, sino también del régimen tecnológico en el que opera cada familia al inicio de la dinámica.


\printbibliography



\end{document}